% ===========================================================================
%  sviluppatore/07-marketplace/03-pagamenti.tex
% ===========================================================================
\section{I pagamenti}\label{sez:pagamenti}

Questa è l'area in cui un errore costa denaro vero, e la regola che la governa
è una sola: \textbf{l'accesso lo concede il webhook, e nessun altro}.

\figuraTikz{pagamento}
  {Il percorso di un acquisto. La freccia sbarrata è quella che esisteva e non
   esiste più.}
  {pagamento}

\subsection{La regola, e perché c'è}

\begin{commento}{dashboard/community.py:1982 $\cdot$ \cod{marketplace\_purchase\_return()}}
Ritorno dal browser dopo il pagamento.

NON concede nulla. Prima questa rotta creava l'acquisto come «paid»
fidandosi della sessione indicata nell'URL: chiunque poteva aprirlo e
ottenere gratis qualunque artefatto. Ora l'accesso lo apre soltanto il
webhook del sito, quando Stripe conferma di aver incassato; qui si
verifica solo se la licenza risulta gia' attiva, per dare all'utente
un messaggio sensato.

Il pagamento e' appena avvenuto e il webhook puo' arrivare un istante
dopo: se la licenza non c'e' ancora non e' un errore, si dice di
attendere.
\end{commento}

\begin{attenzione}{Un identificativo nell'URL non è una prova di pagamento}
È il difetto classico dell'integrazione con un fornitore di pagamenti:
l'utente torna su un indirizzo che contiene l'identificativo della sessione,
e il codice lo tratta come conferma. Ma quell'indirizzo è visibile,
indovinabile e condivisibile. L'unica conferma è quella che arriva
\emph{dal} fornitore, firmata, sul canale del webhook. Qualunque modifica in
quest'area che aggiunga una concessione fuori dal webhook va considerata un
difetto di sicurezza, non una comodità.
\end{attenzione}

\subsection{Quando il webhook non arriva}

I webhook si perdono. Esiste quindi una rotta di riconciliazione, costruita
con lo stesso principio:

\begin{codiceInterno}{python}{Riconciliare senza concedere}{dashboard/community.py:2072}
    def marketplace_riconcilia(artifact_id):
        """Completa un pagamento che il webhook non ha portato a termine.

        Non concede nulla di suo: chiede a Stripe se il pagamento risulta
        incassato, e solo in quel caso fa cio' che avrebbe fatto il webhook.
        """
\end{codiceInterno}

\noindent
«Non concede nulla di suo» è la frase che la rende legittima: non è una
scorciatoia, è una seconda interrogazione alla stessa autorità.

\subsection{Distinguere «non comprato» da «non ho potuto chiedere»}

\begin{commento}{dashboard/community.py:2094 $\cdot$ \cod{marketplace\_perche\_non\_mio()}}
Perche' l'artefatto non risulta acquistato.

Il possesso dipende da due letture: la licenza e l'acquisto. Se
entrambe tornano vuote non si distingue «non l'hai comprato» da «non
sono riuscito a chiedere», e sono cose opposte. Qui si guardano una
per una, dicendo anche se la lettura e' stata rifiutata.
\end{commento}

\noindent
È la stessa distinzione fra zero e «non so» del \vediSez{pagina-nuova},
applicata al caso in cui la differenza è fra «devi pagare» e «c'è un
problema nostro». Senza questa rotta, un utente che ha pagato e non vede
l'artefatto non ha modo di sapere di chi sia il problema, e l'unica mossa
che gli resta è aprire una contestazione.

\subsection{Gli abbonamenti: una rotta sola per due cose}

\begin{commento}{dashboard/community.py:2680 $\cdot$ \cod{abbonamento\_checkout()}}
Apre il pagamento di un piano, o cambia quello in corso.

Una rotta sola per entrambe le cose perche' la domanda dell'utente e'
una sola — «voglio il livello X con cadenza Y» — e il sito sa se
quella richiesta e' un abbonamento nuovo o un cambio con conguaglio.
Due rotte avrebbero costretto l'interfaccia a indovinarlo.
\end{commento}

\noindent
Il corpo della richiesta porta \cod{tier} e \cod{periodo}; le cadenze sono
mensile e annuale, con l'annuale a dieci mensilità (Desk solo annuale). La
disdetta, il cambio di carta e le fatture passano invece dal portale di
fatturazione (\rotta{POST}{/api/abbonamento/portale}), che è una pagina
ospitata dal fornitore: non c'è codice da mantenere per quelle funzioni.

Tre regole stanno sul sito, nella rotta che crea la sessione di pagamento,
e non nell'app:

\begin{itemize}
  \item \textbf{La prova di due mesi.} Se il piano la prevede e l'account non
        ne ha mai avuta una né ha mai avuto un abbonamento su Stripe, la
        sessione porta la fine della prova: mesi di calendario, calcolati
        dal minuto della chiave di idempotenza, così due clic nello stesso
        minuto mandano a Stripe la stessa richiesta. La carta si chiede
        comunque, e il primo addebito parte da solo alla fine della prova. Il
        webhook scrive \cod{subscriptionTrialEnd} finché la prova è in corso
        e \cod{trialUsedAt} la prima volta, che non si svuota più.
  \item \textbf{Desk non si compra dal checkout.} Il piano porta
        \cod{sales\_only} e la rotta risponde \num{409} con il codice
        \cod{sales\_only}, anche a un amministratore.
  \item \textbf{Il piano si riconosce dal codice.} Nei metadati
        dell'abbonamento il checkout scrive il codice del piano accanto al
        livello; il webhook e la rotta delle fatture cercano il piano prima
        dal codice, poi dal prodotto Stripe, poi dal numero. Desk era il
        livello~\num{3} prima che arrivasse Quant, e un abbonamento nato
        allora porta ancora quel numero.
\end{itemize}

\subsection{Le contestazioni}

Un acquirente che ritiene che un artefatto non renda come dichiarato può
aprire una pratica. Due rotte la rendono una verifica e non una discussione:

\begin{center}
\begin{tabular}{@{}L{54mm}L{84mm}@{}}
\toprule
\textbf{Rotta} & \textbf{Che cosa confronta} \\
\midrule
\rigaTenue
\cod{.../verifica-integrita} & I file installati contro i file venduti,
  uno per uno. \\
\cod{.../verifica-performance} & Le metriche dichiarate dal venditore
  contro quelle ottenute rieseguendo qui. \\
\bottomrule
\end{tabular}
\end{center}

\begin{perche}{Perché la verifica è automatica}
Una contestazione fra due persone su che cosa «renda» una strategia non si
risolve discutendo. Le due rotte producono un confronto riproducibile: o i
file sono quelli, o non lo sono; o le metriche si riproducono, o non si
riproducono. L'amministratore che chiude la pratica decide su un risultato,
non su due racconti.
\end{perche}

\subsection{Le variabili d'ambiente}

Le chiavi Stripe, SMTP e simili \emph{non} stanno in un \file{.env} locale:
si gestiscono dal sito e arrivano alla metà locale attraverso il canale~\num{2}
(\vediSez{due-meta}). È il motivo per cui un ambiente di sviluppo può avere
il circuito di pagamento funzionante senza che nessuna chiave stia sul disco
dello sviluppatore.
